%@ pages 411-416
%@ status translated
% Chapter V, section 2: Localisation (pp. 411--416).  Transcribed from the
% page images in reference/source.pdf, not from the OCR text layer, which
% mangles subscripts, primes and the letters $F$/$\mathbf{F}$.
%
% Numbering is automatic (src/preamble.tex).  This section prints
% Proposition 3, Proposition 4 and Proposition 5 (continuing the chapter's
% run, which reaches Proposition 2 on p. 411), Lemma 1 and Lemma 2, one
% corollary which Gabriel prints without a number (hence
% `unnumberedcorollary'), Corollary 1 and Corollary 2, one unnumbered
% definition and Remarks 1--3.
%
% Notation: $F$ is a topologising set of right ideals, $\mathbf{F}$ the
% associated closed subcategory of $\operatorname{mod} A$; $A_d$ is $A$
% viewed as a right module over itself, ${}_{s}B$ is $B$ viewed as a left
% module over itself.  The star conditions $(\star)$ and $(\star\star)$ of
% Proposition 5 are so printed in the original.
%
% Two emendations, both footnoted in the text: in the proof of
% $(a)\Rightarrow(b)$ of Proposition 5 the original prints ``si $a$ est
% nul'' where the argument requires ``si $s.a$ est nul''; in the statement
% of Proposition 5 the original prints ``$A_F$'' for the target of $\psi$,
% where the notation $M_S$ for $M_{F_S}$ just introduced forces $A_S$.

\section{Localisation}

--- Let $A$ be a ring. A topology on $A$ is said to be \emph{linear} (on
the right) if it makes $A$ a topological ring and if there exists a base
of neighbourhoods of $0$ formed of right ideals. In this case, the set
$F$ of open right ideals satisfies the following conditions:

a. If $\mathfrak{m}$ is a right ideal containing an element $\mathfrak{I}$
of $F$, then $\mathfrak{m}$ belongs to $F$.

b. If $\mathfrak{m}\in F$ and if $\mathfrak{I}\in F$, then
$\mathfrak{m}\cap\mathfrak{I}\in F$.

c. If $\mathfrak{I}\in F$ and if $a\in A$, then
$(\mathfrak{I}:a)=\{x\mid a.x\in\mathfrak{I}\}$ belongs to $F$.

We shall say that a set $F$, formed of right ideals of $A$, is
\emph{topologising} if conditions~(a), (b) and~(c) are satisfied. We
shall also say that an $A$-module $M$ is \emph{$F$-negligible} if the
annihilator of every element of $M$ is a right ideal belonging to $F$.
If $\mathbf{F}$ is the full subcategory of $\operatorname{mod} A$ whose
objects are the $F$-negligible $A$-modules, $\mathbf{F}$ is a closed
subcategory (cf.\ Chapter~IV, \S4).

\begin{lemma}
The map $F\mapsto\mathbf{F}$ is a bijection from the set of
topologising sets of right ideals onto the set of closed subcategories
of $\operatorname{mod} A$.
\end{lemma}

If $\mathbf{F}$ is known, one can indeed recover $F$: a right ideal
$\mathfrak{I}$ belongs to $F$ if and only if $A/\mathfrak{I}$ is an
object of $\mathbf{F}$.

If $F$ and $G$ are two topologising sets of right ideals, $F.G$ shall
denote the topologising set associated to the product
$\mathbf{F}.\mathbf{G}$ (Chapter~IV, \S4). It amounts to the same to say
that a right ideal $\mathfrak{m}$ belongs to $F.G$ if, for every
$a\in A$, there exists an element $\mathfrak{I}$ of $F$ such that
$(a.\mathfrak{I}+\mathfrak{m})/\mathfrak{m}$ is $G$-negligible. In
particular, if $\mathfrak{p}$ is an element of $F$ and $\mathfrak{q}$ an
element of $G$, then $\mathfrak{p}.\mathfrak{q}$ belongs to $F.G$.
With these conventions, the closed subcategory $\mathbf{F}$ is
localising if and only if $F.F$ is equal to $\mathbf{F}$. The localising
subcategories of $\operatorname{mod} A$ therefore correspond
bijectively to the sets of right ideals which are topologising and
idempotent.

Let now $F$ be a topologising and idempotent set of right ideals. For
every $A$-module $M$, we choose an $\mathbf{F}$-envelope which we denote
by $u_M:M\to M_F$ (cf.\ Chapter~III). Since $M_F$ is an $A$-module, every
element $m$ of $M_F$ defines an $A$-linear map from $A_d$ into $M_F$:
$a\mapsto m.a$. This map extends in one and only one way to an
$A$-linear map from $A_F=(A_d)_F$ into $M_F$; one thus defines a bilinear
map
\[
M_F\times A_F\to M_F.
\]

We leave to the reader the task of verifying that, if $M$ is equal to
$A_d$, this bilinear map makes $A_F$ a ring. If $M$ is arbitrary, this
map makes $M_F$ an $A_F$-module. In particular, if $\mathfrak{I}$ is a
right ideal of $A$, we may suppose that $\mathfrak{I}_F$ has been
chosen to be a right ideal of $A_F$.

In what follows, we always endow $A_F$ and $M_F$ with the structures we
have just made explicit. We also employ the following notations: $T$ is
the canonical functor from $\operatorname{mod} A$ into
$\operatorname{mod} A/\mathbf{F}$; $S'$ is the functor from
$\operatorname{mod} A/\mathbf{F}$ into $\operatorname{mod} A_F$ induced
by the localisation functor; $\rho$ is the functor which associates to
every $A_F$-module $N$ the underlying $A$-module (for the $A$-module
structure defined by $u_A$). With these notations, $\rho\circ S'$ is a
functor adjoint to $T$.

\begin{lemma}
Let $M$ be an $A$-module and $N$ an $A_F$-module. The map
$\rho(N,M_F)$ from $\operatorname{Hom}_{A_F}(N,M_F)$ into
$\operatorname{Hom}_A(\rho N,\rho M_F)$ is bijective.
\end{lemma}

It is indeed clear that $\rho(N,M_F)$ is an injection. Let $f$ therefore
be an $A$-linear map from $N$ into $M_F$. If $n$ is an element of $N$,
we want to show that $f(n.a)=f(n).a$ for every $a$ in $A_F$. But the
maps $g:a\mapsto f(n.a)$ and $h:a\mapsto f(n).a$ are $A$-linear, and one
has the equality $g\circ u_A=h\circ u_A$. We deduce the equality of $g$
and of $h$ (definition of $\mathbf{F}$-closed objects, Lemma~1, \S2,
Chapter~III).\hfill Q.E.D.

Let $F'$ be the set of right ideals $\mathfrak{I}$ of $A_F$ such that
$\rho(A_F/\mathfrak{I})$ is $F$-negligible. It is clear that $F'$ is
topologising and idempotent, and that an $A_F$-module $N$ is
$F'$-negligible if and only if $\rho N$ is $F$-negligible. We can
therefore speak of the canonical functor $T'$ from
$\operatorname{mod} A_F$ into the quotient category
$\operatorname{mod} A_F/\mathbf{F}'$.

\begin{proposition}
Let $A$ be a ring and $F$ a topologising and idempotent set of right
ideals. With the notations above, the functor $T\circ\rho$ defines by
passage to the quotient an equivalence between
$\operatorname{mod} A_F/\mathbf{F}'$ and
$\operatorname{mod} A/\mathbf{F}$.
\end{proposition}

Let indeed $\mathbf{B}$ be the quotient category
$\operatorname{mod} A/\mathbf{F}$; let us show that the functor $S'$ is
adjoint to $T\circ\rho$: if $M$ is an object of $\mathbf{B}$ and $N$ an
object of $\operatorname{mod} A_F$, the abelian group
$\operatorname{Hom}_{\mathbf{B}}(T\rho N,M)$ is identified with
$\operatorname{Hom}_A(\rho N,\rho S'M)$, since $\rho\circ S'$ is adjoint
to $T$. By Lemma~2, $\operatorname{Hom}_A(\rho N,\rho S'M)$ ``is
identified'' with $\operatorname{Hom}_{A_F}(N,S'M)$. This proves that
$S'$ is adjoint to $T\circ\rho$. The proposition therefore follows from
Proposition~5 (Chapter~III).

\begin{unnumberedcorollary}
An $A_F$-module $N$ is $\mathbf{F}'$-closed (resp.\ $\mathbf{F}'$-closed
and injective) if and only if $\rho N$ is $\mathbf{F}$-closed (resp.\
$\mathbf{F}$-closed and injective).
\end{unnumberedcorollary}

This corollary follows from the preceding proposition and from the
corollary to Proposition~3 (Chapter~III).

If $M$ is an $A$-module, we denote by $\mathbf{F}M$ the largest
$F$-negligible submodule of $M$. The abelian group underlying $M_F$ is
identified with $\operatorname{Hom}_A(A_d,M_F)$; this last group is
itself isomorphic to
$\operatorname{Hom}_{\operatorname{mod} A/\mathbf{F}}(A_d,M)$, that is,
to the inductive limit
$\varinjlim_{I\in F}\operatorname{Hom}_A(I,M/\mathbf{F}M)$.

It is easy to endow this inductive limit with a structure of $A$-module
such that the preceding isomorphism from $M_F$ onto
$\varinjlim_{I\in F}\operatorname{Hom}_A(I,M/\mathbf{F}M)$ is an
isomorphism of $A$-modules. The proposition which follows is a
consequence of this:

\begin{proposition}
Let $A$ be a ring and $F$ a topologising and idempotent set of right
ideals. The localisation functor defined by $F$ is isomorphic to the
functor
$M\leadsto\varinjlim_{I\in F}\operatorname{Hom}_A(I,M/\mathbf{F}M)$.
\end{proposition}

\begin{corollary}
The hypotheses are those of the proposition. One supposes moreover that
the injective envelope of an object of $\mathbf{F}$ belongs to
$\mathbf{F}$. The localisation functor defined by $F$ is then
isomorphic to the functor
$M\leadsto\varinjlim_{I\in F}\operatorname{Hom}_A(I,M)$.
\end{corollary}

Consider indeed the following exact sequence:
\[
0\to\varinjlim_{I\in F}\operatorname{Hom}_A(I,\mathbf{F}M)
\xrightarrow{i}\varinjlim_{I\in F}\operatorname{Hom}_A(I,M)
\xrightarrow{p}\varinjlim_{I\in F}\operatorname{Hom}_A(I,M/\mathbf{F}M)
\xrightarrow{d}\varinjlim_{I\in F}\operatorname{Ext}^1_A(I,\mathbf{F}M).
\]
One checks at once that
$\varinjlim_{I\in F}\operatorname{Hom}_A(I,N)$ is zero if $N$ is
$F$-negligible. It follows that $p$ is an injection. On the other hand,
the functor
$N\leadsto\varinjlim_{I\in F}\operatorname{Ext}^1_A(I,N)$ is the
satellite of the functor
$N\leadsto\varinjlim_{I\in F}\operatorname{Hom}_A(I,N)$. If $N$ is
$F$-negligible, there exists an exact sequence
\[
0\to N\xrightarrow{u} I\xrightarrow{v} N'\to 0,
\]
where $I$ is an $F$-negligible injective $A$-module. One deduces that
$\varinjlim_{I\in F}\operatorname{Ext}^1_A(I,N)$ is zero if $N$ is
$F$-negligible. In particular, the map $d$ is surjective.

\begin{corollary}
Let $A$ be a ring and $F$ a topologising and idempotent set of right
ideals. Suppose that there exists a cofinal subset $G$ of $F$, formed
of finitely generated right ideals. If the localisation functor defined
by $F$ is exact, this functor is isomorphic to the functor
$M\leadsto M\otimes_A A_F$; moreover, the categories
$\operatorname{mod} A/\mathbf{F}$ and $\operatorname{mod} A_F$ are then
equivalent.
\end{corollary}

Let us first show that the functor
$M\leadsto\varinjlim_{I\in G}\operatorname{Hom}_A(I,M/\mathbf{F}M)$
commutes with direct sums. This is obviously true for the functors
$M\leadsto\mathbf{F}M$ and $M\leadsto M/\mathbf{F}M$; if $N$ is a
finitely generated $A$-module, the functor $\operatorname{Hom}_A(N,.)$
commutes with direct sums. It follows that the functor
$M\leadsto\operatorname{Hom}_A(I,M/\mathbf{F}M)$ commutes with direct
sums if $I$ is an element of $G$. The first assertion of the corollary
follows (see Proposition~1\textit{bis} and the remark preceding it).

Let furthermore $\mathbf{B}$ be the quotient category
$\operatorname{mod} A/\mathbf{F}$, let $T$ be the canonical functor from
$\operatorname{mod} A$ into $\mathbf{B}$ and let $S$ be a functor
adjoint to $T$. The functor $\operatorname{Hom}_{\mathbf{B}}(A_F,.)$ is
isomorphic to the functor $\operatorname{Hom}_A(A,S.)$; since $S$ is
exact, $A_F$ is a projective object of $\mathbf{B}$; since $S$ commutes
with direct sums, the same holds for
$\operatorname{Hom}_{\mathbf{B}}(A_F,.)$. The last assertion therefore
follows from Corollary~1 to Proposition~1.

Here is an application of Corollary~2: let $S$ be a multiplicative set
of $A$ (a subset of $A$ such that $a.b\in S$ if $a\in S$ and $b\in S$).
The datum of $S$ makes it possible to define a topologising and
idempotent set $F_S$ of right ideals: a right ideal $\mathfrak{I}$
belongs to $F_S$ if, for every $a\in A$, there exists an $s\in S$ such
that $a.s$ belongs to $\mathfrak{I}$. It amounts to the same to say
that an $A$-module $M$ is $F_S$-negligible if every element $m$ of $M$
is annihilated by an element $s$ of $S$.

Every ideal belonging to $F_S$ obviously meets $S$; the converse is
true if $\forall s\in S$, $\forall a\in A$, $\exists t\in S$,
$\exists b\in A$ such that $a.t=s.b$. We shall recover this condition;
for this, we write $M_S$ instead of $M_{F_S}$:

\begin{definition}
Let $\varphi:A\to B$ be a homomorphism of rings with unit element and
$S$ a multiplicative set of $A$. One says that $(B,\varphi)$ is a
\emph{right ring of fractions} of $A$ for $S$ if:
\begin{description}
\item[\normalfont a.] $\operatorname{Ker}\varphi$ is $F_S$-negligible;
in other words, if $\varphi(a)$ is zero, there exists $s\in S$ such that
$a.s=0$.
\item[\normalfont b.] The image under $\varphi$ of every element of $S$
is invertible (on the left and on the right).
\item[\normalfont c.] Every element $b$ of $B$ is of the form
$b=\varphi(a).\varphi(s)^{-1}$, $a\in A$, $s\in S$.
\end{description}
\end{definition}

\begin{proposition}
Let $A$ be a ring and $S$ a multiplicative set of $A$. The following
assertions are equivalent:
\begin{description}
\item[\normalfont a.] There exists a right ring of fractions of $A$ for
$S$.
\item[\normalfont b.] $S$ satisfies the following conditions:
\end{description}
$(\star)\quad\forall s\in S,\ \forall a\in A,\ \exists t\in S,\
\exists b\in A$ such that $a.t=s.b$.

$(\star\star)\quad$ if $a\in A$, if $s\in S$ and if $s.a=0$, there
exists $t\in S$ such that $a.t=0$.
\begin{description}
\item[\normalfont c.] The image of an element of $S$ under the map
$u_A:A\to A_S$ is invertible.
\end{description}
If these conditions are satisfied and if $(B,\varphi)$ is a right ring
of fractions of $A$ for $S$, there exists one and only one isomorphism
$\psi$ from $B$ onto $A_S$ such that
$\psi\circ\varphi=u_A$.\ednote{The original prints $A_F$ here; with the
notation $M_S$ for $M_{F_S}$ introduced above, the target can only be
$A_S$.} The functor $M\leadsto M_S$ is then exact and is isomorphic to
the functor $M\leadsto M\otimes_A A_S$.
\end{proposition}

$(a)\Rightarrow(b)$: assertion $(\star\star)$ is indeed clear, for
$\varphi(a)$ is zero if $s.a$ is zero.\ednote{The original prints ``si
$a$ est nul''; since $\varphi(s)$ is invertible,
$\varphi(s.a)=0$ gives $\varphi(a)=0$, and it is $s.a$ which must be
zero for the conclusion $a.t=0$ to follow from~(a) of the definition.}
On the other hand, $\varphi(s)^{-1}.\varphi(a)$ is an element of $B$,
and is therefore of the form
\[
\varphi(s)^{-1}.\varphi(a)=\varphi(c).\varphi(r)^{-1},\qquad r\in S,
\qquad c\in A.
\]
One deduces the equality
$\varphi(a).\varphi(r)=\varphi(s).\varphi(c)$. In other words,
$a.r-s.c$ belongs to $\operatorname{Ker}\varphi$; there consequently
exists a $u\in S$ such that $a.r.u-s.c.u=0$. One shall put $t=r.u$ and
$b=c.u$.

$(b)\Rightarrow(c)$: we first show that conditions $(\star)$ and
$(\star\star)$ imply the exactness of the functor $M\leadsto M_S$. Let
us give ourselves a diagram
\[
\begin{tikzcd}
& I \arrow[d, "u"] \\
M \arrow[r, "p"] & N \arrow[r] & 0.
\end{tikzcd}
\]
where $I$ belongs to $F_S$ and where $M$ and $\operatorname{Ker}p$ are
$F_S$-closed. It suffices to show that $I$ contains an ideal
$\mathfrak{m}\in F_S$ such that the restriction of $u$ to
$\mathfrak{m}$ is of the form $p\circ v$, where $v:\mathfrak{m}\to M$
[Proposition~7~(b), Chapter~III].

If $s$ belongs to $I\cap S$, $u(s)$ is of the form $p(m)$, where
$m\in M$; moreover, $m.a$ is zero if $s.a$ is zero. One shall therefore
put $\mathfrak{m}=s.A$ and $v(s.a)=m.a$.

The exactness of the functor $M\leadsto M_S$ implies that $u_A(s)$ is
invertible on the right for every $s\in S$ (let indeed $g_s$ be the map
$a\mapsto s.a$; since $\operatorname{Coker}g_s$ is $F_S$-negligible,
$g_s$ induces an epimorphism from $A_S$ onto $A_S$). It follows that
$u_A(s).(s'.u_A(s)-1)$ is zero if $s'$ is such that
$u_A(s).s'=1$. Since $u_A(s)$ is an element regular on the right and on
the left [condition $(\star\star)$], one has $s'.u_A(s)=1$.

$(c)\Rightarrow(a)$: this is clear, as is the uniqueness of the ring of
fractions. The last assertion follows from Corollary~2 to
Proposition~4.

\begin{remark}
Conditions $(\star)$ and $(\star\star)$ are always satisfied when $S$ is
contained in the centre of $A$. One then recovers localisation with
respect to the centre (Chapter~III, \S5).
\end{remark}

\begin{remark}
If conditions $(\star)$ and $(\star\star)$ of Proposition~5 are
satisfied, one can exhibit a construction of $M_S$ closer to
established traditions: for this, one agrees to say that two elements
$(m,s)$ and $(n,t)$ of $M\times S$ are equivalent if there exist two
elements $u$ and $v$ of $A$ such that $s.u=t.v\in S$ and $m.u=n.v$.
The set underlying $M_S$ is then chosen equal to the quotient of
$M\times S$ by the equivalence relation we have just defined. The
definition of the laws of addition and multiplication is easy.
\end{remark}

\begin{remark}
We leave to the reader the task of proving that, if $A$ is commutative,
$A_F$ is commutative for every topologising and idempotent set of ideals
$F$; moreover, the map $\mathfrak{p}\mapsto\mathfrak{p}_F$ is a
bijection from the set of prime ideals of $A$ not belonging to $F$ onto
the set of prime ideals of $A_F$ not belonging to $F'$.
\end{remark}
